
\chapter{大容量存储系统}

\section{磁盘结构与访问时间}
\concept{磁盘结构}
\begin{points}
  \item 磁盘由一个或多个盘片组成，盘片表面覆盖磁性材料。
  \item 每个盘面对应一个磁头，磁臂沿半径方向移动。
  \item 磁道是盘面上的同心圆，扇区是磁道上的固定大小区域。
  \item 柱面是所有盘面上半径相同的磁道集合。
  \item 一个物理块可用柱面号、磁头号、扇区号表示。
\end{points}
\plain{这部分先建立“盘片-盘面-磁头-磁道-扇区-柱面”的空间结构图，后面访问时间、调度算法、坏块管理都建立在这个结构上。}
\exam{概念题常考柱面、磁道、扇区的区别；计算题则常默认你已经能据此理解磁头为什么要移动。}


\concept{磁盘访问时间}
\begin{points}
  \item 寻道时间：磁臂移动到目标柱面/磁道所需时间。
  \item 旋转延迟：目标扇区旋转到磁头下方所需时间，平均为旋转一周时间的一半。
  \item 传输时间：数据从磁盘读出或写入磁盘的时间。
  \item 总访问时间约为：寻道时间 + 旋转延迟 + 传输时间。
  \item 机械磁盘优化重点通常是减少寻道时间。
\end{points}
\plain{磁盘访问时间要分解着看，机械盘里通常最贵的是寻道和旋转等待，不是数据真正传过去那一下。}
\exam{PPT 还给了完整 I/O 时间计算：定位时间 + 传输时间 + 控制器负载时间，做题别漏最后一项。}


\concept{逻辑块映射与磁盘连接方式}
\begin{points}
  \item 文件系统常把磁盘看成逻辑块的一维数组，逻辑块按一定规则顺序映射到实际扇区。
  \item 逻辑到物理的地址转换并不总是简单线性，因为可能存在坏扇区、不同磁道扇区数不完全一致等情况。
  \item 存储连接方式常见有 DAS、NAS、SAN：分别对应主机直接连接、经网络文件访问、经专用存储网络访问。
  \item 总线、主机控制器、磁盘控制器共同参与主机与磁盘的通信。
\end{points}
\plain{题目写“逻辑块号”时，不等于它天然就是磁盘上的物理位置编号；中间往往还隔着控制器和映射规则。}
\exam{PPT 单独讲了 DAS/NAS/SAN，选择题里常考三者“直接接主机 / 通过网络文件协议 / 通过专用存储网络”的差别。}


\concept{磁盘、磁带与 SSD}
\begin{points}
  \item 磁带是顺序访问设备，容量大、成本低，适合备份和归档，但随机访问极慢。
  \item 磁盘是直接访问设备，支持随机定位块，是典型通用二级存储。
  \item SSD 基于闪存，读取快、抗震、噪音低，但价格较高且擦写寿命有限。
\end{points}
\plain{这三类设备的关键区别是访问方式和适用场景，不是背一堆参数。}
\exam{常考顺序存取设备和直接存取设备代表，或比较 HDD、SSD、磁带的优缺点。}


\section{磁盘调度算法}
\concept{FCFS}
\begin{points}
  \item 按磁盘请求到达顺序服务。
  \item 优点：公平、简单。
  \item 缺点：磁头来回移动可能很大，平均寻道距离长。
\end{points}
\plain{FCFS 在磁盘里就是不做任何优化地按到达顺序走，所以磁头可能来回折返很多次。}
\exam{常考按给定请求序列计算 FCFS 的总移动磁道数。}


\concept{SSTF 最短寻道时间优先}
\begin{points}
  \item 每次选择距离当前磁头位置最近的请求。
  \item 优点：平均寻道距离通常较短。
  \item 缺点：远处请求可能长期等待，产生饥饿。
\end{points}
\plain{SSTF 每次都挑最近的请求，平均移动少，但远端请求可能一直被近处新请求压着，出现饥饿。}
\exam{常考按请求序列逐步选择“离当前磁头最近”的请求，并分析饥饿问题。}


\concept{SCAN 电梯算法}
\begin{points}
  \item 磁头像电梯一样沿一个方向移动，服务沿途请求，到达端点后反向。
  \item 优点：比 FCFS 稳定，避免频繁来回。
  \item 缺点：靠近两端的请求等待时间可能不同。
\end{points}
\plain{SCAN 的关键不是“最近优先”，而是“沿一个方向扫过去再回头”，这样能避免磁头反复折返。}
\exam{有些教材把只扫到最远请求就掉头的版本也直接叫 SCAN/C-SCAN；考试前要按老师 PPT 的定义作答。}


\concept{C-SCAN 循环扫描}
\begin{points}
  \item 磁头只沿一个方向服务，到达一端后快速返回另一端，不在返回途中服务。
  \item 提供更均匀的等待时间。
  \item 适合请求分布较均匀的系统。
\end{points}
\plain{C-SCAN 像循环单向电梯：服务时始终朝一个方向走，所以两端请求受到的对待更平均。}
\exam{算总移动距离时，别忘了它从一端“跳回”另一端那段距离通常也要计入。}


\concept{LOOK 与 C-LOOK}
\begin{points}
  \item LOOK 不一定走到磁盘端点，只走到当前方向最远请求处就反向。
  \item C-LOOK 只在有请求的范围内循环扫描，从一端最远请求跳到另一端最远请求。
  \item 它们是 SCAN/C-SCAN 的实际优化版本，减少无意义移动。
\end{points}
\plain{LOOK/C-LOOK 的改进点就是“不走冤枉路”：没有请求的尾部磁道就不必特意扫到头。}
\exam{老师 PPT 明说它们是 SCAN/C-SCAN 的改良版，若题目要求综合性能，常优先想到 LOOK 类。}


\concept{N-Step-SCAN}
\begin{points}
  \item 把磁盘请求队列按长度 $N$ 划分成若干批次。
  \item 批次之间按 FCFS 处理，每个批次内部再按 SCAN 调度。
  \item 它可以减轻普通 SCAN 的磁臂粘着问题，避免新到请求不断打断当前扫描。
  \item 当 $N$ 很大时接近普通 SCAN；当 $N=1$ 时退化为 FCFS。
\end{points}
\plain{核心思想是“先把请求分批，再对每批做电梯调度”，这样新请求不会一直把磁头吸住。}
\exam{若题目给出 $N$ 和新到请求，要先切批，再分别算每一批内部的扫描路径。}


\concept{FSCAN}
\begin{points}
  \item FSCAN 是 N-Step-SCAN 的简化版，只维护两个队列。
  \item 当前扫描只处理旧队列；扫描期间新到达的请求进入新队列。
  \item 当前队列处理完后，再把新队列作为下一轮扫描对象。
  \item 这样可避免新请求持续涌入导致当前扫描迟迟结束不了。
\end{points}
\plain{可以把 FSCAN 看成“双缓冲版 SCAN”：一边扫当前队列，一边收集下一轮请求。}
\exam{常和普通 SCAN 对比考，问它为什么更能避免磁臂粘着。}


\concept{如何选择磁盘调度算法}
\begin{points}
  \item 磁盘调度算法效果受请求数量、请求分布、文件分配方式等因素共同影响。
  \item 不存在对所有场景都绝对最优的算法，因此系统常保留可替换的调度策略。
  \item SSTF 和 LOOK 往往是较常见的默认选择。
  \item SCAN 与 C-SCAN 在高负载场景下更有利于避免饥饿、维持较稳定性能。
\end{points}
\plain{这类题不要死背“哪一个最好”，真正结论是没有通杀算法，只能按负载特征取舍。}
\exam{简答题若问“为什么没有最优磁盘调度算法”，核心要点是请求分布和系统负载是动态变化的。}


\section{磁盘存储优化分布}
\concept{记录排列会影响 I/O 总时间}
\begin{points}
  \item 在旋转型磁盘上，记录在磁道上的排列顺序会影响旋转等待时间。
  \item 如果逻辑上相邻的记录物理上也紧挨着，但处理上一条记录需要时间，则下一条记录可能已经转过磁头。
  \item 因此记录的物理分布需要结合“读出时间 + 处理时间”做错位安排。
  \item 这一优化减少的是旋转延迟，不是寻道时间。
\end{points}
\plain{上一条记录刚处理完时，下一条记录最好正好转到磁头下，而不是刚错过去。}
\exam{常考比较普通顺序分布和优化分布下的总处理时间。}


\concept{优化分布的直觉}
\begin{points}
  \item 若 CPU 处理当前记录要花若干毫秒，就应让下一条逻辑记录在若干扇区之后出现。
  \item 这样 CPU 处理结束时，下一条逻辑记录刚好到达磁头位置。
  \item 题目若给出“一圈时间、每条记录读出时间、处理时间、每道记录数”，通常按转过多少个扇区后恰好接上下一次读取来排布。
  \item SSD 不依赖机械旋转，因此不存在这种扇区优化分布问题。
\end{points}
\plain{它像排队卡节奏：不是让人全挤在窗口前，而是每次轮到时下一个人刚好走到窗口。}
\exam{常考计算题或简答：为什么优化排布能把总时间从“多等整圈”降到“几乎无额外等待”。}


\section{磁盘管理与容错}
\concept{磁盘格式化}
\begin{points}
  \item 低级格式化把磁盘划分为扇区，为每个扇区建立控制信息。
  \item 分区把磁盘划分为一个或多个逻辑区域。
  \item 高级格式化/逻辑格式化在分区上建立文件系统数据结构，如空闲空间表、根目录、inode 区等。
\end{points}
\plain{这三步是从“原始盘面”到“能存文件的文件系统分区”的构建过程：先切扇区，再切分区，再建文件系统。}
\exam{常考低级格式化、逻辑分区、高级格式化三者分别解决什么问题，别混成一步。}


\concept{扇区、块与簇}
\begin{points}
  \item 扇区 sector 是磁盘的物理存储单位。
  \item 块 block 是操作系统进行磁盘 I/O 时常用的逻辑传输单位，通常由多个扇区组成。
  \item 簇 cluster 常作为文件系统分配文件空间的单位，也可能由多个扇区或多个块组成。
  \item 做题时要先分清题目说的是物理扇区、I/O 块还是文件分配簇。
\end{points}
\plain{扇区是硬盘原生小格子，块和簇是操作系统为传输和分配方便再打的包。}
\exam{常考“磁盘 I/O 以什么为单位”“文件系统分配空间以什么为单位”这类辨析。}


\concept{原始磁盘 raw disk}
\begin{points}
  \item 原始磁盘或原始分区指不建立普通文件系统、直接由应用管理块的磁盘区域。
  \item 它减少了文件系统带来的额外管理开销，适合数据库等追求高性能的场景。
  \item 代价是通用性差，很多管理责任要交给应用自己处理。
\end{points}
\plain{raw disk 就是“别走文件系统这层，我直接控制这块盘”。}
\exam{常考 raw disk 为什么可能更快，以及它牺牲了哪些便利性。}


\concept{引导块与坏块}
\begin{points}
  \item 引导块保存启动操作系统所需的引导代码或引导信息。
  \item 坏块是无法可靠读写的磁盘块。
  \item 坏块管理可通过扇区备用、坏块链表、控制器重映射等方式处理。
\end{points}
\plain{磁盘管理不仅是读写，还包括初始化、启动、坏块处理和容错。}
\exam{PPT 还给了引导流程示例：ROM/BIOS -> MBR -> 分区引导扇区 -> 引导管理器 -> 内核，知道层次即可。}


\concept{坏块管理与重映射}
\begin{points}
  \item 发现坏块后，系统可把它记录到坏块表或坏块链表中，后续分配时直接跳过。
  \item 现代磁盘还常预留备用块或备用扇区，用于对坏块做透明重映射。
  \item 对上层来说，目标是屏蔽物理缺陷，尽量不让文件系统继续碰到坏块。
\end{points}
\plain{坏块一般不是去“修”，而是尽快标记并绕开。}
\exam{常考坏块链表、备用块、控制器重映射各自起什么作用。}


\concept{交换空间管理}
\begin{points}
  \item 交换空间用于保存被换出的进程映像或页面。
  \item 交换空间可作为专门分区，也可作为文件系统中的交换文件。
  \item 专门分区管理简单且效率高；交换文件更灵活但可能受文件系统开销影响。
\end{points}
\plain{交换空间不是普通存文件的地方，而是虚拟内存把页面或进程映像临时换出去的后备区域。}
\exam{常考交换空间的用途，以及它和普通文件存储区域的区别。}


\concept{交换文件与交换分区比较}
\begin{points}
  \item 交换文件实现简单，配置灵活，能直接建在现有文件系统里。
  \item 但访问交换文件通常还要经过文件系统的相关结构，可能带来额外开销。
  \item 专用交换分区不依赖普通文件系统，访问路径更直接，更便于专门优化。
  \item 因为交换访问很频繁，系统常更偏向专用交换分区。
\end{points}
\plain{交换区是虚拟内存的后备仓库，用得勤，所以越少绕文件系统这层越好。}
\exam{常考为什么独立交换分区通常比交换文件效率更高。}


\concept{RAID 磁盘容错}
\begin{points}
  \item RAID 通过磁盘冗余提高性能和可靠性。
  \item RAID 0 条带化提高性能但无冗余。
  \item RAID 1 镜像保存完整副本，可靠性高但空间开销大。
  \item RAID 5 使用分布式奇偶校验，兼顾容量和容错。
  \item RAID 6 可容忍两块磁盘故障，可靠性更高。
\end{points}
\plain{RAID 要分清“提速”和“容错”不是总能同时最大化；不同级别是在容量、性能、可靠性间做取舍。}
\exam{除了 0/1/5/6，PPT 还出现了 RAID 2/3/4：它们分别引入海明码校验、专用奇偶校验盘、块级奇偶校验盘，虽少用但可能出选择题。}


\concept{并行交叉存取与 RAID 2/3/4}
\begin{points}
  \item 并行交叉存取把一个盘块拆成多个子块并分散到多个磁盘上并行读写，以缩短传输时间。
  \item RAID 2 以位/字节级拆分并配合海明码校验，结构复杂，商业上较少使用。
  \item RAID 3 使用一块专用奇偶校验盘，适合大块顺序传输场景，如科学计算、图像处理。
  \item RAID 4 与 RAID 3 类似，但按块级拆分数据，仍使用专用奇偶校验盘。
\end{points}
\plain{RAID 2/3/4 的共同思路都是“条带化 + 冗余校验”，差别主要在数据拆分粒度和校验组织方式。}
\exam{若题目问“哪几级依赖专用校验盘”，要想到 RAID 3/4；若问“分布式奇偶校验”，那是 RAID 5。}


\concept{RAID 5}
\begin{points}
  \item RAID 5 采用块级条带化，并把奇偶校验信息分布到各个磁盘上，而不是集中到一块专用校验盘。
  \item 每个驱动器通常都有独立的数据通道，可独立进行读写，适合 I/O 繁忙的事务处理场景。
  \item 它兼顾了容量利用率、读写性能和容错能力，是工程上较常见的一类 RAID 方案。
  \item 与 RAID 3/4 相比，RAID 5 的关键改进是消除了专用校验盘的单点瓶颈。
\end{points}
\plain{RAID 5 的记忆重点不是“也有校验”，而是“校验分散到所有盘上”，所以既能容错，又不容易把一块校验盘压成热点。}
\exam{选择题和比较题常把 RAID 3/4/5 放一起考：3/4 有专用校验盘，5 是分布式奇偶校验。}


\concept{RAID 6 与 RAID 7}
\begin{points}
  \item RAID 6 在冗余校验方面比 RAID 5 更强，可进一步容忍更多故障，可靠性更高。
  \item 课件表述里，RAID 6 设置了一个可快速访问的异步校验盘，并具有独立数据访问通道，性能优于部分早期奇偶校验方案。
  \item RAID 7 在阵列控制层集成了更智能的实时控制软件，可独立于主机运行，尽量减少主机 CPU 负担。
  \item 记忆时抓住趋势：级别越高，通常越强调控制智能化和更强的故障容忍能力。
\end{points}
\plain{这两级不用背工程细节，考试里更重要的是知道它们是在 RAID 5 之后继续往“更强容错/更强控制能力”方向演化。}
\exam{若选择题专门问 RAID 6/7，抓“RAID 6 更强容错，RAID 7 更智能、对主机依赖更少”通常就够了。}


\concept{磁盘容错层次}
\begin{points}
  \item 低级容错主要针对磁盘表面缺陷，例如双份目录、双份文件分配表、写后读校验、热修复重定向。
  \item 中级容错主要针对驱动器或控制器故障，例如磁盘镜像和磁盘双工。
  \item 更高层容错关注系统级持续服务能力。
\end{points}
\plain{容错不只是一句 RAID，它还分“防单个坏块”“防整块盘坏”“防整套系统失效”几个层次。}
\exam{常考双份目录、写后读校验属于哪一级容错，以及镜像和双工分别冗余了什么。}

\concept{双份目录和双份文件分配表}
\begin{points}
  \item 目录和文件分配表都是文件系统运行所依赖的关键数据结构，一旦损坏会直接影响文件定位与空间管理。
  \item 为提高可靠性，可在不同磁盘或同一磁盘的不同区域保存两份：一份主目录/主文件分配表，一份备份目录/备份文件分配表。
  \item 当主副本损坏时，系统可切换到备份副本继续工作。
  \item 系统启动时通常还要检查两份数据结构的一致性。
\end{points}
\plain{这类方法保护的不是“普通文件内容”，而是文件系统自己的核心账本；账本一坏，整套文件管理都会受影响。}
\exam{简答题若问低级容错的典型措施，写出“双份目录、双份文件分配表、主备切换、一致性检查”会比只写“做备份”更稳。}

\concept{热修复重定向与写后读校验}
\begin{points}
  \item \term{写后读校验}：每次把数据块写入磁盘后，立即再读出并与原写入数据比较，确认写入结果是否正确。
  \item 若校验失败且重写后仍不一致，就可判定该盘块存在缺陷。
  \item \term{热修复重定向}：把待写数据改写到预留的热修复重定向区，并登记映射关系，供以后访问时查找。
  \item 这类方法适合缺陷数量较少时对局部坏块做补救，而不必立刻放弃整块磁盘。
\end{points}
\plain{写后读校验负责“发现写坏了没有”，热修复重定向负责“既然这块不可靠，就把数据挪到备用区并记好新地址”。}
\exam{选择题常考两者分工：前者偏检测，后者偏绕开坏块并继续提供可访问位置。}


\concept{磁盘镜像与磁盘双工}
\begin{points}
  \item \term{磁盘镜像}：两台磁盘驱动器连接到同一个磁盘控制器，主盘写入后再把同样数据写到备份盘，主盘故障时启用备份盘。
  \item 磁盘镜像主要冗余的是\key{磁盘驱动器}本身，但若共享的那一个控制器失效，两块盘仍可能一起不可用。
  \item \term{磁盘双工}：两台磁盘驱动器分别接到两个不同的磁盘控制器，同时保存完全相同的数据副本。
  \item 磁盘双工比镜像更进一步，因为它同时冗余了\key{磁盘驱动器和磁盘控制器}，可靠性更高，但成本也更高。
\end{points}
\plain{镜像和双工最稳的区分方式不是背图，而是看“两个备份盘是不是还共用同一个控制器”：共用的是镜像，不共用的是双工。}
\exam{比较题常直接问两者区别，标准答法就是“镜像冗余盘，双工冗余盘加控制器”，并顺带说明双工可靠性更高。}


\concept{第三级系统容错}
\begin{points}
  \item 第三级系统容错建立在前两级容错基础上，目标是不因单台服务器故障而中断整体服务。
  \item 典型做法是主从服务器镜像：主服务器正常对外服务，从服务器实时保存主服务器的关键内存和磁盘数据副本。
  \item 当主服务器故障时，从服务器接替成为新的服务节点；故障排除后，两台服务器再重新同步。
  \item 这一级关注点已不只是“某块盘坏了怎么办”，而是“整个系统如何持续提供服务”。
\end{points}
\plain{前两级主要还在“保住磁盘和控制器”，第三级已经上升到“即使整台主服务器挂了，业务也别停”。}
\exam{PPT 单独讲了第三级容错，简答题常考它和前两级的层次区别：前者偏系统接管，后者偏磁盘/控制器冗余。}
